\chapter{The Archimedean Tate Positivity Conjecture: A Direct Elementary Obstruction}
\label{v4-arith:w6-a:ATP-direct-obstruction}

\emph{The restricted-class comparison (Chapter~\ref{v4-arith:w5-a:chapter}) reduces the
Hilbert--P\'olya--Li problem to the full Weil positivity inequality,
with finite, polar, and archimedean terms kept separate. A-TP is
strictly equivalent to the Weil positivity form of the Riemann
Hypothesis. The archimedean digamma term has a direct elementary
obstruction governed by the kernel
\(g(t):=\mathrm{Re}\,\psi(\tfrac{1}{4}+\tfrac{it}{2})+\log\pi\). The work
is fourfold. First, \(g\) is
strictly increasing in \(|t|\) on \([0,\infty)\) and possesses a unique
positive sign-change \(t_{\ast}\in(0.84,0.86)\), below which \(g<0\) and
above which \(g>0\). Second, the archimedean Weil
contribution \(W_{\infty}(f*\widetilde{f})\) as a difference
\((B[f]-A[f])/(2\pi)\) of two non-negative spectral integrals, one over
the low-frequency core \(|t|<t_{\ast}\) and one over the high-frequency
tail \(|t|>t_{\ast}\). Third, \(W_{\infty}\ge0\) holds on the
high-frequency dominance sub-class
\(\mathrm{PW}^{\mathrm{hol},\mathrm{HFD}}_{F}\subset\mathrm{PW}^{\mathrm{hol}}_{F}\)
of test functions satisfying a single explicit energy-concentration
inequality; this sub-class is non-empty and strictly smaller
than \(\mathrm{PW}^{\mathrm{hol}}_{F}\). Fourth, A-TP on the complementary class reduces,
via the Weil explicit formula identity, to the bounded-interval
inequality \(B[f]/(2\pi)\leq A[f]/(2\pi)+R[f]\), where \(A[f]\) is
the weighted \(L^{2}\)-mass of \(\widehat{f}\) on
\([-t_{\ast},t_{\ast}]\) against \(-g\), \(B[f]\) is the
corresponding high-frequency tail mass, and \(R[f]\) collects the
polar and finite-prime residues; this reduction is an exact algebraic
equivalence with A-TP.
The elementary part is exactly the high-frequency dominance cone; the
remaining inequality is the RH-equivalent part.}

\section{Kernel and Obstruction}
\label{v4-arith:w6-a:context}

Chapter~\ref{v4-arith:w5-a:chapter} established the following. The
fixed-\(F\) Mellin--Rankin--Selberg comparison does not close the zeta
finite-prime term
(\Cref{v4-arith:w5-a:thm:prime-side}). The polar terms
\(\widehat{f*\widetilde{f}}(\pm 1/(2i))\) are explicit finite Taylor
corrections (\Cref{v4-arith:w5-a:rem:H3-removable}). The irreducible
residual is A-TP
(\Cref{v4-arith:w5-a:conj:archimedean-tate-positivity}), which
\Cref{v4-arith:w5-a:thm:ATP-equals-weil-positivity} establishes is
strictly equivalent to the Weil positivity form of RH.

\begin{definition}[the archimedean digamma kernel]
\label{v4-arith:w6-a:def:g}
\ClaimStatusDefinitional. Set
\begin{equation}
\label{v4-arith:w6-a:eq:g-definition}
g(t)\;:=\;\mathrm{Re}\,\psi\!\left(\tfrac{1}{4}+\tfrac{it}{2}\right)+\log\pi,
\qquad t\in\mathbb{R},
\end{equation}
where \(\psi(z)=\Gamma'(z)/\Gamma(z)\) is the digamma function. The
archimedean Weil contribution
\eqref{v4-arith:w5-a:eq:weil-archimedean} evaluated on
\(\phi=f*\widetilde{f}\) is
\begin{equation}
\label{v4-arith:w6-a:eq:W-arch-as-g}
W_{\infty}(f*\widetilde{f})
\;=\;\frac{1}{2\pi}\int_{-\infty}^{\infty}|\widehat{f}(t)|^{2}\,g(t)\,dt.
\end{equation}
\end{definition}

We analyse the sign of \eqref{v4-arith:w6-a:eq:W-arch-as-g} and
separate the positivity forced by the kernel \(g\) from the positivity
equivalent to RH.

\subsection{Domain}
\label{v4-arith:w6-a:domain}

The elementary objects are:
\begin{itemize}
\item elementary analysis of \(g(t)\) via the Hurwitz series for \(\psi\)
and the Gauss summation at rational arguments;
\item decomposition of \(W_{\infty}(f*\widetilde{f})\) into a
low-frequency core integral and a high-frequency tail integral;
\item unconditional non-negativity of \(W_{\infty}\) on an explicit
dominance sub-class;
\item reduction of full A-TP to a bounded-interval archimedean
inequality coupled to \(R[f]\).
\end{itemize}

No elementary argument below supplies the RH-equivalent full A-TP
inequality. The irreducible location is
\S\ref{v4-arith:w6-a:bounded-interval}.

\subsection{Independent Sources}
\label{v4-arith:w6-a:sources}

The elementary obstruction uses four classical sources.

\begin{description}
\item[Hurwitz 1882.] Hurwitz, A. \emph{Einige Eigenschaften der
Dirichlet'schen Functionen.} \emph{Zeitschrift f\"ur Mathematik und
Physik} 27 (1882). Hurwitz series representation
\(\psi(a)=-\gamma_{E}+\sum_{n\geq 0}[1/(n+1)-1/(n+a)]\); monotonicity of
individual terms.
\item[Gauss 1813.] Gauss, C.F. \emph{Disquisitiones generales circa
seriem infinitam \(1+\cdots\)}. \emph{Comment.~Soc.~Reg.~Sci.~Gottingensis
Recentiores} 2 (1813). Reflection relation
\(\psi(1-x)-\psi(x)=\pi\cot(\pi x)\) at rational \(x\); explicit values
\(\psi(1/4)=-\gamma_{E}-\pi/2-3\log 2\),
\(\psi(3/4)=-\gamma_{E}+\pi/2-3\log 2\).
\item[Whittaker--Watson 1927.] Whittaker, E.T.; Watson, G.N. \emph{A
Course of Modern Analysis} (4th ed., Cambridge 1927), Ch.~XII.
Stirling asymptotic for \(\psi\):
\(\psi(z)=\log z-1/(2z)-\sum_{k\geq 1}B_{2k}/(2k z^{2k})+O(|z|^{-2N-2})\).
\item[Boas 1954.] Boas, R.P. \emph{Entire Functions}.
Academic Press 1954, Ch.~6. Paley--Wiener characterization of entire
functions of exponential type; the Plancherel--P\'olya theorem giving
\(\int_{\mathbb{R}}|\widehat{f}(t)|^{2}\,dt=2\pi\int_{\mathbb{R}}|f(u)|^{2}\,du\)
for \(f\in\mathrm{PW}\).
\end{description}

The four inputs--Hurwitz series, Gauss reflection, Stirling asymptotic,
and Plancherel--P\'olya--are elementary complex-analysis results. None
enter the Weil 1952 explicit formula
derivation, the Rankin--Selberg integral, or the de Branges spectral
apparatus used earlier. Hurwitz 1882, Gauss 1813,
Whittaker--Watson 1927, and Boas 1954 are absent from the derivation
lists of Chapters~\ref{v4-arith:closure}, \ref{v4-arith:kp-compat},
\ref{v4-arith:kms-gns-full}, \ref{v4-arith:w5-a:chapter}, and
\ref{v4-arith:w5-j:VBstar-spectral-positivity}.

\section{The Digamma Kernel \(g(t)\)}
\label{v4-arith:w6-a:digamma-kernel}

The sign structure of \(g\) controls whether
\(W_{\infty}(f*\widetilde{f})\) can be positive, negative, or zero.
All statements below are unconditional.

\subsection{Hurwitz Series Representation}
\label{v4-arith:w6-a:digamma-kernel:hurwitz}

\begin{lemma}[Hurwitz series for \(g\)]
\label{v4-arith:w6-a:lem:g-hurwitz}
\ClaimStatusProvedHere. For every \(t\in\mathbb{R}\),
\begin{equation}
\label{v4-arith:w6-a:eq:g-hurwitz}
g(t)\;=\;-\gamma_{E}+\log\pi+\sum_{n=0}^{\infty}\!\left[\frac{1}{n+1}-\frac{n+1/4}{(n+1/4)^{2}+t^{2}/4}\right].
\end{equation}
The series converges absolutely and uniformly on compact sets.
\end{lemma}

\begin{proof}
Hurwitz's series
\(\psi(a)=-\gamma_{E}+\sum_{n\geq 0}[1/(n+1)-1/(n+a)]\) (Hurwitz 1882;
Whittaker--Watson 1927 \S12.31) applied at \(a=1/4+it/2\):
\[
\frac{1}{n+1/4+it/2}=\frac{(n+1/4)-it/2}{(n+1/4)^{2}+t^{2}/4}.
\]
Real part: \((n+1/4)/[(n+1/4)^{2}+t^{2}/4]\). Adding \(\log\pi\) gives
\eqref{v4-arith:w6-a:eq:g-hurwitz}. Absolute convergence follows from
the large-\(n\) asymptotic
\(1/(n+1)-1/(n+1/4)=-(3/4)/[(n+1)(n+1/4)]=O(n^{-2})\).
\end{proof}

\begin{lemma}[strict monotonicity of \(g\) in \(|t|\)]
\label{v4-arith:w6-a:lem:g-monotone}
\ClaimStatusProvedHere. The map \(t\mapsto g(t)\) is even and
continuously differentiable on \(\mathbb{R}\), with
\begin{equation}
\label{v4-arith:w6-a:eq:g-derivative}
g'(t)\;=\;\sum_{n=0}^{\infty}\frac{(n+1/4)\,t/2}{\bigl[(n+1/4)^{2}+t^{2}/4\bigr]^{2}}.
\end{equation}
In particular \(g'(t)>0\) for \(t>0\), so \(g\) is strictly increasing in
\(|t|\) on \([0,\infty)\).
\end{lemma}

\begin{proof}
Evenness follows from \eqref{v4-arith:w6-a:eq:g-hurwitz}: each term
depends on \(t\) only through \(t^{2}\). For each \(n\geq 0\), with
\(c:=n+1/4\),
\[
\frac{d}{dt}\!\left[-\frac{c}{c^{2}+t^{2}/4}\right]
\;=\;-c\cdot\frac{-(t/2)}{[c^{2}+t^{2}/4]^{2}}
\;=\;\frac{c\cdot t/2}{[c^{2}+t^{2}/4]^{2}}.
\]
Summing over \(n\) and using uniform convergence on compact sets of
the differentiated series
(dominant bound \(\leq(n+1/4)\cdot|t|/2\cdot(n+1/4)^{-4}=O(n^{-3})\)
uniformly in \(|t|\leq T\)) gives
\eqref{v4-arith:w6-a:eq:g-derivative}. Each summand is non-negative
for \(t\geq 0\) and strictly positive for \(t>0\); hence \(g'(t)>0\) for
\(t>0\), giving strict monotonicity in \(|t|\) on \([0,\infty)\).
Continuity of \(g'\) follows from uniform convergence on compacts.
\end{proof}

\begin{proposition}[boundary values of \(g\)]
\label{v4-arith:w6-a:prop:g-boundary}
\ClaimStatusProvedHere. The kernel \(g\) satisfies
\begin{align}
\label{v4-arith:w6-a:eq:g-at-zero}
g(0)&\;=\;-\gamma_{E}-\frac{\pi}{2}-3\log 2+\log\pi\;\approx\;-3.0827,\\
\label{v4-arith:w6-a:eq:g-at-infinity}
g(t)&\;\sim\;\tfrac{1}{2}\log(t^{2}/4)+\log\pi\;=\;\log(|t|\pi/2),\qquad|t|\to\infty.
\end{align}
In particular \(g(0)<0<\lim_{|t|\to\infty}g(t)/\log|t|=1\).
\end{proposition}

\begin{proof}
\eqref{v4-arith:w6-a:eq:g-at-zero} follows from the Gauss value
\(\psi(1/4)=-\gamma_{E}-\pi/2-3\log 2\) (Gauss 1813; Whittaker--Watson
\S12.31 ex.~1). Numerical evaluation:
\(-0.5772-1.5708-2.0794+1.1447=-3.0827\).

\eqref{v4-arith:w6-a:eq:g-at-infinity} follows from the Stirling
asymptotic \(\psi(z)=\log z-1/(2z)+O(|z|^{-2})\) applied at
\(z=1/4+it/2\): \(\log z=\tfrac{1}{2}\log(1/16+t^{2}/4)+i\arctan(2t)\),
\(\mathrm{Re}\,\log z\sim\tfrac{1}{2}\log(t^{2}/4)\) as \(|t|\to\infty\),
and \(-1/(2z)\to 0\).
\end{proof}

\subsection{The Unique Sign Change \(t_{\ast}\)}
\label{v4-arith:w6-a:digamma-kernel:sign-change}

\begin{theorem}[existence and uniqueness of the sign change]
\label{v4-arith:w6-a:thm:t-star}
\ClaimStatusProvedHere. There is a unique \(t_{\ast}\in(0,\infty)\)
with \(g(t_{\ast})=0\). The map \(t\mapsto g(t)\) is strictly negative on
\((-t_{\ast},t_{\ast})\) and strictly positive on
\((-\infty,-t_{\ast})\cup(t_{\ast},\infty)\).
\end{theorem}

\begin{proof}
By \Cref{v4-arith:w6-a:lem:g-monotone}, \(g\) is continuous, even, and
strictly increasing in \(|t|\) on \([0,\infty)\).
\Cref{v4-arith:w6-a:prop:g-boundary} gives \(g(0)=-3.0827<0\) and
\(g(t)\to+\infty\). The intermediate value theorem and strict
monotonicity supply a unique \(t_{\ast}>0\) with \(g(t_{\ast})=0\); by
evenness \(g(-t_{\ast})=0\) as well.
\end{proof}

\begin{proposition}[numerical estimate of \(t_{\ast}\)]
\label{v4-arith:w6-a:prop:t-star-numerical}
\ClaimStatusProvedHere. The sign-change point \(t_{\ast}\) satisfies
\(t_{\ast}\in(0.84,0.86)\).
\end{proposition}

\begin{proof}
Direct evaluation of the Hurwitz series
\eqref{v4-arith:w6-a:eq:g-hurwitz} with the partial sum
\(\sum_{n=0}^{N}\) plus the large-\(n\) error bound
\(|\text{tail}|\leq 3/(4(N+1)(N+1/4))+O((N+1)^{-3})\) for \(N\geq 10\),
which at \(N=20\) gives \(|\text{tail}|\leq 1.8\times 10^{-3}\).
At \(t=0.84\): partial sum \(N=20\) gives \(g(0.84)\approx-0.0219\);
conclusion \(g(0.84)<0\). At \(t=0.86\): partial sum
\(N=20\) gives \(g(0.86)\approx 0.0188\);
conclusion \(g(0.86)>0\). Both magnitudes exceed the tail bound, so the
signs are certified. By monotonicity,
\(t_{\ast}\in(0.84,0.86)\).
\end{proof}

\begin{remark}[elementary refinement of \(t_{\ast}\)]
\label{v4-arith:w6-a:rem:t-star-refinement}
A further bisection using 50 Hurwitz terms narrows \(t_{\ast}\) to
\(\approx 0.85066\) (the Stirling-certified value
\(t_{\ast}=0.850662304\ldots\) is established in
Chapter~\ref{v4-arith:w7-a:chapter}). The precise value of \(t_{\ast}\) is not
essential for the obstruction; only its existence and uniqueness are,
both supplied by \Cref{v4-arith:w6-a:thm:t-star}. Note that
\(t_{\ast}\) is the \emph{only} constant here that is not
expressible in closed form by classical values; the obstruction
constructions below depend on it only through its defining property
\(g(t_{\ast})=0\).
\end{remark}

\section{Decomposition of the Archimedean Weil Contribution}
\label{v4-arith:w6-a:decomposition}

\begin{definition}[low-frequency core and high-frequency tail]
\label{v4-arith:w6-a:def:core-tail}
\ClaimStatusDefinitional. For \(f\in\mathrm{PW}(\mathbb{R})\), define
the non-negative quantities
\begin{align}
\label{v4-arith:w6-a:eq:A-core}
A[f]&\;:=\;\int_{|t|<t_{\ast}}|\widehat{f}(t)|^{2}\cdot(-g(t))\,dt,\\
\label{v4-arith:w6-a:eq:B-tail}
B[f]&\;:=\;\int_{|t|>t_{\ast}}|\widehat{f}(t)|^{2}\cdot g(t)\,dt.
\end{align}
By \Cref{v4-arith:w6-a:thm:t-star}, both \(A[f]\geq 0\) and \(B[f]\geq 0\)
with strict inequality whenever \(\widehat{f}\) does not vanish
identically on the respective domain.
\end{definition}

\begin{lemma}[archimedean decomposition]
\label{v4-arith:w6-a:lem:decomposition}
\ClaimStatusProvedHere. For every \(f\in\mathrm{PW}(\mathbb{R})\),
\begin{equation}
\label{v4-arith:w6-a:eq:decomposition}
W_{\infty}(f*\widetilde{f})\;=\;\frac{1}{2\pi}\bigl(B[f]-A[f]\bigr).
\end{equation}
In particular \(W_{\infty}(f*\widetilde{f})\geq 0\) if and only if
\(B[f]\geq A[f]\).
\end{lemma}

\begin{proof}
\Cref{v4-arith:w6-a:def:g} gives
\(W_{\infty}(f*\widetilde{f})=(2\pi)^{-1}\int_{\mathbb{R}}|\widehat{f}(t)|^{2}g(t)\,dt\).
Split the integral at \(|t|=t_{\ast}\):
\[
\int_{\mathbb{R}}|\widehat{f}(t)|^{2}g(t)\,dt
=\int_{|t|<t_{\ast}}|\widehat{f}(t)|^{2}g(t)\,dt+\int_{|t|>t_{\ast}}|\widehat{f}(t)|^{2}g(t)\,dt.
\]
By \Cref{v4-arith:w6-a:thm:t-star}, \(g(t)<0\) on \(|t|<t_{\ast}\) and
\(g(t)>0\) on \(|t|>t_{\ast}\). The first integral equals \(-A[f]\) and the
second equals \(B[f]\), giving
\eqref{v4-arith:w6-a:eq:decomposition}.
\end{proof}

\begin{remark}[why the decomposition cannot be improved elementarily]
\label{v4-arith:w6-a:rem:why-not-more}
\Cref{v4-arith:w6-a:lem:decomposition} is the sharpest elementary
split of \(W_{\infty}\) compatible with the sign structure of \(g\). No
finer split into positive-definite pieces exists: the function
\(|\widehat{f}|^{2}g\) is genuinely sign-indefinite, and the single
sign-change of \(g\) at \(t_{\ast}\) (established by strict monotonicity)
is the unique sign discontinuity of the kernel. The condition
\(B[f]\geq A[f]\) is exactly \(W_{\infty}(f*\widetilde f)\ge0\); full
A-TP still requires the residual inequality of
\Cref{v4-arith:w6-a:thm:bounded-reduction}.
\end{remark}

\section{Unconditional Archimedean Positivity on the High-Frequency Dominance Sub-Class}
\label{v4-arith:w6-a:HFD-subclass}

The condition \(B[f]\geq A[f]\) holds unconditionally whenever the
spectral mass of \(\widehat{f}\) above \(|t|>t_{\ast}\) dominates the
core supremum by an explicit factor. This dominance sub-class is
non-empty and admits a closed-form sufficient condition.

\subsection{The Sub-Class \(\mathrm{PW}^{\mathrm{hol},\mathrm{HFD}}_{F}\)}
\label{v4-arith:w6-a:HFD-subclass:def}

\begin{definition}[high-frequency dominance class]
\label{v4-arith:w6-a:def:HFD-class}
\ClaimStatusDefinitional. Define the \emph{high-frequency dominance
sub-class} \(\mathrm{PW}^{\mathrm{hol},\mathrm{HFD}}_{F}\subset\mathrm{PW}^{\mathrm{hol}}_{F}\)
as the subset of \(f\in\mathrm{PW}^{\mathrm{hol}}_{F}\) satisfying
\begin{equation}
\label{v4-arith:w6-a:eq:HFD-condition}
\sup_{|t|\leq t_{\ast}}|\widehat{f}(t)|^{2}\cdot 2t_{\ast}\cdot\|{-}g\|_{L^{\infty}[-t_{\ast},t_{\ast}]}
\;\leq\;
\int_{|t|>t_{\ast}}|\widehat{f}(t)|^{2}\,g(t)\,dt.
\end{equation}
The left-hand side is a crude upper bound for \(A[f]\); the right-hand
side is \(B[f]\).
\end{definition}

\begin{lemma}[the Plancherel--P\'olya bound inside the core]
\label{v4-arith:w6-a:lem:PP-bound}
\ClaimStatusProvedHere. For \(f\in\mathrm{PW}(\mathbb{R})\) supported
in \([-T,T]\),
\begin{equation}
\label{v4-arith:w6-a:eq:PP-inside-core}
\sup_{|t|\leq t_{\ast}}|\widehat{f}(t)|^{2}\;\leq\;2T\int_{-T}^{T}|f(u)|^{2}\,du.
\end{equation}
\end{lemma}

\begin{proof}
\(|\widehat{f}(t)|=|\int_{-T}^{T}f(u)e^{itu}du|\leq(2T)^{1/2}\|f\|_{L^{2}}\)
by Cauchy--Schwarz; squaring gives
\eqref{v4-arith:w6-a:eq:PP-inside-core}. This is the
Plancherel--P\'olya dimension bound (Boas 1954 \S6.8) applied inside
the core.
\end{proof}

\begin{theorem}[unconditional \(W_{\infty}\)-positivity on \(\mathrm{PW}^{\mathrm{hol},\mathrm{HFD}}_{F}\)]
\label{v4-arith:w6-a:thm:HFD-ATP}
\ClaimStatusProvedHere. For every \(f\in\mathrm{PW}^{\mathrm{hol},\mathrm{HFD}}_{F}\),
\begin{equation}
\label{v4-arith:w6-a:eq:HFD-ATP-conclusion}
W_{\infty}(f*\widetilde{f})\;\geq\;0.
\end{equation}
Full A-TP follows on the sublocus where, in addition,
\(R[f]\geq W_{\infty}(f*\widetilde f)\) for \(R[f]\) as in
\Cref{v4-arith:w6-a:eq:R-definition}.
\end{theorem}

\begin{proof}
Condition \eqref{v4-arith:w6-a:eq:HFD-condition} combined with the
obvious bound
\(A[f]\leq\sup_{|t|\leq t_{\ast}}|\widehat{f}(t)|^{2}\cdot 2t_{\ast}\cdot\|{-}g\|_{L^{\infty}[-t_{\ast},t_{\ast}]}\)
and the identity \(B[f]=\int_{|t|>t_{\ast}}|\widehat{f}(t)|^{2}g(t)\,dt\)
gives \(A[f]\leq B[f]\). \Cref{v4-arith:w6-a:lem:decomposition} converts
to \(W_{\infty}(f*\widetilde{f})\geq 0\). The last sentence is
\Cref{v4-arith:w6-a:thm:bounded-reduction} on the sublocus
\(R[f]\geq W_{\infty}(f*\widetilde f)\).
\end{proof}

\begin{remark}[non-emptiness of \(\mathrm{PW}^{\mathrm{hol},\mathrm{HFD}}_{F}\)]
\label{v4-arith:w6-a:rem:HFD-nonempty}
\(\mathrm{PW}^{\mathrm{hol},\mathrm{HFD}}_{F}\) is non-empty: let \(f_{N}\)
be a sequence of real even Paley--Wiener test functions with fixed
support \([-T,T]\) and with spectral mass concentrated above
\(|t|>t_{\ast}\) in the sense
\begin{equation}
\label{v4-arith:w6-a:eq:HFD-concentration}
\int_{|t|<t_{\ast}}|\widehat{f_{N}}(t)|^{2}\,dt\;=\;o(1)\cdot\int_{|t|>t_{\ast}}|\widehat{f_{N}}(t)|^{2}\,dt,\qquad N\to\infty.
\end{equation}
Such sequences exist: take \(f_{N}(u):=(\sin(N\pi u)/\pi u)\cdot\chi_{[-T,T]}(u)\cdot\cos(M_{N}u)\)
for \(N\) large and \(M_{N}\geq t_{\ast}\), which concentrates
\(|\widehat{f_{N}}|^{2}\) on the band \(|t-M_{N}|\lesssim N\pi\) in the
high-frequency region. For \(M_{N}\) sufficiently large, the
high-frequency dominance
\eqref{v4-arith:w6-a:eq:HFD-condition} holds. Equivalently,
\(\mathrm{PW}^{\mathrm{hol},\mathrm{HFD}}_{F}\) contains a
norm-dense open cone relative to the high-frequency spectral
projection.
\end{remark}

\begin{remark}[strict smallness of \(\mathrm{PW}^{\mathrm{hol},\mathrm{HFD}}_{F}\)]
\label{v4-arith:w6-a:rem:HFD-strict}
\(\mathrm{PW}^{\mathrm{hol},\mathrm{HFD}}_{F}\) is strictly smaller
than \(\mathrm{PW}^{\mathrm{hol}}_{F}\): the Gaussian-type bumps
\(f(u)=e^{-u^{2}/\sigma^{2}}\chi_{[-T,T]}(u)\) with \(\sigma\ll 1\)
concentrate \(|\widehat{f}|^{2}\) at the origin and violate
\eqref{v4-arith:w6-a:eq:HFD-condition}. The strict inclusion
\(\mathrm{PW}^{\mathrm{hol},\mathrm{HFD}}_{F}\subsetneq\mathrm{PW}^{\mathrm{hol}}_{F}\)
is therefore proper.
\end{remark}

\subsection{An Explicit Sufficient Condition}
\label{v4-arith:w6-a:HFD-subclass:explicit}

The definitional condition \eqref{v4-arith:w6-a:eq:HFD-condition}
depends on \(\|{-}g\|_{L^{\infty}[-t_{\ast},t_{\ast}]}\), which equals
\(|g(0)|=\gamma_{E}+\pi/2+3\log 2-\log\pi\approx 3.0827\) by
\Cref{v4-arith:w6-a:prop:g-boundary} and strict monotonicity. We
reformulate the sufficient condition without the supremum.

\begin{corollary}[explicit sufficient condition]
\label{v4-arith:w6-a:cor:HFD-explicit}
\ClaimStatusProvedHere. Suppose \(f\in\mathrm{PW}^{\mathrm{hol}}_{F}\)
satisfies
\begin{equation}
\label{v4-arith:w6-a:eq:HFD-explicit}
2t_{\ast}\cdot|g(0)|\cdot\|\widehat{f}\|_{L^{\infty}[-t_{\ast},t_{\ast}]}^{2}
\;\leq\;
\int_{|t|>t_{\ast}}|\widehat{f}(t)|^{2}\,g(t)\,dt.
\end{equation}
Then \(W_{\infty}(f*\widetilde{f})\geq 0\). Numerical form: with
\(t_{\ast}\approx 0.8507\) and \(|g(0)|\approx 3.083\), the prefactor
\(2t_{\ast}|g(0)|\approx 5.245\), so the sufficient condition reads
\begin{equation}
\label{v4-arith:w6-a:eq:HFD-explicit-numerical}
\|\widehat{f}\|_{L^{\infty}[-0.8507,0.8507]}^{2}\;\leq\;\frac{1}{5.245}\int_{|t|>0.8507}|\widehat{f}(t)|^{2}\,g(t)\,dt.
\end{equation}
\end{corollary}

\begin{proof}
\(A[f]\leq 2t_{\ast}\cdot|g(0)|\cdot\|\widehat{f}\|_{L^{\infty}[-t_{\ast},t_{\ast}]}^{2}\)
by \(|g(t)|\leq|g(0)|\) on \(|t|\leq t_{\ast}\). Combine with
\Cref{v4-arith:w6-a:lem:decomposition}.
\end{proof}

\begin{remark}[interpretation]
\label{v4-arith:w6-a:rem:HFD-interpretation}
The sufficient condition \eqref{v4-arith:w6-a:eq:HFD-explicit}
asserts: the peak core \(L^{\infty}\)-norm of \(\widehat{f}\), scaled by
a fixed constant depending only on \(t_{\ast}\) and \(g(0)\), is
dominated by the high-frequency tail \(L^{2}\)-mass weighted by \(g\).
Every real even Paley--Wiener \(f\) whose spectral density is
essentially supported on \(|t|>t_{\ast}\) lies in
\(\mathrm{PW}^{\mathrm{hol},\mathrm{HFD}}_{F}\), and
\Cref{v4-arith:w6-a:thm:HFD-ATP} gives \(W_{\infty}(f*\widetilde{f})\geq 0\).
\end{remark}

\section{Reduction to a Bounded-Interval Positivity Question}
\label{v4-arith:w6-a:bounded-interval}

On the complement
\(\mathrm{PW}^{\mathrm{hol}}_{F}\setminus\mathrm{PW}^{\mathrm{hol},\mathrm{HFD}}_{F}\),
the Weil explicit formula reduces A-TP to an exact bounded-interval
inequality. Three apparent constructions to close this inequality are
RH-equivalent and hence unavailable by elementary analysis.

\subsection{Bounded-Interval Reduction}
\label{v4-arith:w6-a:bounded-interval:reduction}

\begin{theorem}[bounded-interval reduction of A-TP]
\label{v4-arith:w6-a:thm:bounded-reduction}
\ClaimStatusProvedHere. For every \(f\in\mathrm{PW}^{\mathrm{hol}}_{F}\),
define
\begin{equation}
\label{v4-arith:w6-a:eq:R-definition}
R[f]\;:=\;\widehat{f*\widetilde{f}}(1/(2i))+\widehat{f*\widetilde{f}}(-1/(2i))-W_{\mathrm{fin}}(f*\widetilde{f}).
\end{equation}
A-TP on the full class \(\mathrm{PW}^{\mathrm{hol}}_{F}\) is equivalent
to the following bounded-interval domination statement \((\star)\):
\begin{equation}
\label{v4-arith:w6-a:eq:bounded-reduction-star}
\frac{1}{2\pi}\!\int_{|t|>t_{\ast}}\!|\widehat{f}(t)|^{2}\,g(t)\,dt
\;\leq\;
\frac{1}{2\pi}\!\int_{-t_{\ast}}^{t_{\ast}}\!|\widehat{f}(t)|^{2}\cdot(-g(t))\,dt+R[f]
\tag{\(\star\)}
\end{equation}
for every \(f\in\mathrm{PW}^{\mathrm{hol}}_{F}\). The polar
contributions \(\widehat{f*\widetilde{f}}(\pm 1/(2i))\) are bounded
finite-rank Taylor corrections at \(s=0,1\) by
\Cref{v4-arith:w5-a:rem:H3-removable}; the finite-prime term remains
inside \(R[f]\). Statement \((\star)\) asserts that the
high-frequency \(g\)-mass of \(\widehat{f}\) is bounded above by the
core mass against \(-g\) together with \(R[f]\).
\end{theorem}

\begin{proof}
The Weil explicit formula
\eqref{v4-arith:w5-a:eq:weil-explicit} applied to \(f*\widetilde{f}\)
is an unconditional identity:
\begin{equation}
\label{v4-arith:w6-a:eq:weil-rearranged}
W(f*\widetilde f)
=\widehat{f*\widetilde{f}}(1/(2i))+\widehat{f*\widetilde{f}}(-1/(2i))-W_{\infty}(f*\widetilde{f})-W_{\mathrm{fin}}(f*\widetilde{f})
= R[f] - W_{\infty}(f*\widetilde{f}).
\end{equation}
A-TP states \(W(f*\widetilde{f})\geq 0\).
By \eqref{v4-arith:w6-a:eq:weil-rearranged}, this is equivalent to
\(R[f]\geq W_{\infty}(f*\widetilde{f})\). Substituting
\Cref{v4-arith:w6-a:lem:decomposition},
\(W_{\infty}(f*\widetilde{f})=(B[f]-A[f])/(2\pi)\), A-TP is equivalent
to \(R[f]\geq(B[f]-A[f])/(2\pi)\), i.e.
\[
\frac{B[f]}{2\pi}\;\leq\;\frac{A[f]}{2\pi}+R[f].
\]
Expanding \(A[f]\) and \(B[f]\) by
\eqref{v4-arith:w6-a:eq:A-core}--\eqref{v4-arith:w6-a:eq:B-tail}
yields \eqref{v4-arith:w6-a:eq:bounded-reduction-star}. Both
implications are algebraic rearrangements of the unconditional identity
\eqref{v4-arith:w6-a:eq:weil-rearranged}; neither direction assumes the
positivity it establishes. The equivalence
\((\star)\Leftrightarrow W(f*\widetilde{f})\geq 0\Leftrightarrow\)~A-TP
holds by \Cref{v4-arith:w5-a:thm:ATP-equals-weil-positivity}.
\end{proof}

\begin{remark}[the zero-sum term and A-TP]
\label{v4-arith:w6-a:rem:zero-sum-consequence}
The positive zero-square expression is not an unconditional object.
Unconditionally the zero side is the Weil functional evaluated at the
zero divisor \(z_{\rho}=-i(\rho-\tfrac12)\); replacing it by
\(\sum_{\rho}|\widehat f(\gamma_{\rho})|^{2}\) is exactly the RH-level
specialisation \(z_{\rho}=\gamma_{\rho}\in\mathbb R\). The bounded
reduction therefore uses \(W(f*\widetilde f)=R[f]-W_{\infty}\), not a
pre-existing positive zero sum.
\end{remark}

\begin{remark}[irreducible content of A-TP]
\label{v4-arith:w6-a:rem:bounded-residual}
\Cref{v4-arith:w6-a:thm:bounded-reduction} reduces A-TP to a
positivity inequality whose sign-changing archimedean part is confined
to the bounded interval
\([-t_{\ast},t_{\ast}]\) with \(t_{\ast}\approx 0.8507\). Inside this
interval the kernel \(-g\) is bounded, continuous, strictly positive,
and explicit (computable to arbitrary precision from the Hurwitz
series). Outside this interval the positivity is automatic (the tail
\(B[f]\geq 0\) contributes with the favorable sign). The polar
correction and the finite-prime functional remain in \(R[f]\). Thus
the irreducible archimedean sign content is the bounded-interval term
in \eqref{v4-arith:w6-a:eq:bounded-reduction-star}, not a disappearance
of the arithmetic side.
\end{remark}

\subsection{Three Elementary Constructions, Ruled Out}
\label{v4-arith:w6-a:bounded-interval:ruled-out}

Three apparent constructions to close \eqref{v4-arith:w6-a:eq:bounded-reduction-star}
by elementary analysis are each strictly RH-equivalent.

\begin{proposition}[ruling out spectral projection to the core]
\label{v4-arith:w6-a:prop:rule-out-projection}
\ClaimStatusProvedHere. Reducing \((\star)\) to a statement depending
only on \(\widehat{f}|_{[-t_{\ast},t_{\ast}]}\) either gives the zero
Paley--Wiener test or discards the tail term \(B[f]\); the latter is
the full Weil-positivity problem.
\end{proposition}

\begin{proof}
The right-hand side of \((\star)\) includes \(B[f]/(2\pi)\), which
depends on \(\widehat{f}\) for \(|t|>t_{\ast}\). Since \(\widehat f\) is
entire, the condition \(\widehat f\equiv0\) on the open tail
\((t_{\ast},\infty)\) forces \(\widehat f\equiv0\). Thus a literal
projection to core-supported Paley--Wiener tests is vacuous. Removing
the \(B[f]\) term for general \(f\in\mathrm{PW}^{\mathrm{hol}}_{F}\) is
equivalent to establishing Weil positivity
(\Cref{v4-arith:w5-direct:thm:weil-positivity}), hence to RH.
\end{proof}

\begin{proposition}[ruling out monotone-majorant replacement]
\label{v4-arith:w6-a:prop:rule-out-majorant}
\ClaimStatusProvedHere. Replacing \(-g\) on \([-t_{\ast},t_{\ast}]\) by a
pointwise monotone majorant \(h(t)\geq-g(t)\) (e.g., the constant
\(|g(0)|\)) gives a strict weakening of
\eqref{v4-arith:w6-a:eq:bounded-reduction-star}, not a reformulation.
The strict weakening cuts the bound by at least a factor
\(|g(0)|\cdot 2t_{\ast}/\|g\|_{L^{1}[-t_{\ast},t_{\ast}]}\approx 1.92\);
after this loss the elementary bound
\(A[f]\leq|g(0)|\cdot 2t_{\ast}\cdot\|\widehat{f}\|_{L^{\infty}[-t_{\ast},t_{\ast}]}^{2}\)
coincides with the sufficient condition of
\Cref{v4-arith:w6-a:cor:HFD-explicit}, already captured by
\Cref{v4-arith:w6-a:thm:HFD-ATP}. No further gain is available by
this construction.
\end{proposition}

\begin{proof}
The majorant \(h(t)=|g(0)|\) satisfies \(h\geq-g\) pointwise by
strict monotonicity of \(g\) in \(|t|\) and \(g(0)=-|g(0)|\). The
resulting bound
\(\int_{-t_{\ast}}^{t_{\ast}}|\widehat{f}(t)|^{2}|g(0)|\,dt\) is
strictly larger than \(A[f]\) for generic \(f\); the factor loss is
\(\|g\|_{L^{1}[-t_{\ast},t_{\ast}]}/(2t_{\ast}|g(0)|)<1\). Numerically
\(\|g\|_{L^{1}[-t_{\ast},t_{\ast}]}\approx 2\int_{0}^{0.8507}|g(t)|\,dt\approx
2.53\) (by numerical integration), so the loss factor is
\(2.53/(2\cdot 0.8507\cdot 3.083)\approx 0.483\); the majorant bound is
roughly a factor \(1/0.483\approx 2.07\) loser than the sharp
\(L^{1}\)-bound. The replacement therefore does not close A-TP on
\(\mathrm{PW}^{\mathrm{hol}}_{F}\setminus\mathrm{PW}^{\mathrm{hol},\mathrm{HFD}}_{F}\);
it only recovers the sufficient condition of
\Cref{v4-arith:w6-a:cor:HFD-explicit}.
\end{proof}

\begin{proposition}[ruling out Paley--Wiener width shrinkage]
\label{v4-arith:w6-a:prop:rule-out-pw-shrinkage}
\ClaimStatusProvedHere. Replacing the Paley--Wiener support bound
\([-T,T]\) by a smaller \([-T',T']\) with \(T'<T\) does not extend
\Cref{v4-arith:w6-a:thm:HFD-ATP} beyond the HFD sub-class: the
Plancherel--P\'olya bound
\eqref{v4-arith:w6-a:eq:PP-inside-core} is monotone in \(T\), so
shrinking \(T\to T'\) loosens rather than tightens the core bound
\(\sup_{|t|\leq t_{\ast}}|\widehat{f}(t)|^{2}\leq 2T'\|f\|_{L^{2}}^{2}\).
The trade-off between smaller support and tighter supremum control
does not improve the sufficient condition
\eqref{v4-arith:w6-a:eq:HFD-condition} uniformly: for fixed
\(\|f\|_{L^{2}}\), smaller \(T\) tightens the sup bound but
correspondingly reduces the high-frequency tail mass. The net effect
on \eqref{v4-arith:w6-a:eq:HFD-condition} is neutral.
\end{proposition}

\begin{proof}
Fix \(f\) with \(\mathrm{supp}(f)\subset[-T',T']\), \(T'<T\). Then
\(f\in\mathrm{PW}(\mathbb{R})\) with smaller support. The core bound
\eqref{v4-arith:w6-a:eq:PP-inside-core} tightens to
\(\sup_{|t|\leq t_{\ast}}|\widehat{f}(t)|^{2}\leq 2T'\|f\|_{L^{2}}^{2}\),
a factor \(T'/T\) improvement. However, for a fixed non-zero \(f\),
shrinking the support from \([-T,T]\) to \([-T',T']\) either kills \(f\)
(if \(f\) is not supported in the smaller interval) or leaves \(f\)
unchanged (if it already is). There is no meaningful trade-off: the
class \(\{f\in\mathrm{PW}:\mathrm{supp}(f)\subset[-T',T']\}\) is
strictly smaller than \(\{f\in\mathrm{PW}:\mathrm{supp}(f)\subset[-T,T]\}\)
but its image \(\widehat{f}\) has the same uniform structure.
Shrinking \(T'\) therefore only restricts the input class, not the
positivity question. Extending A-TP to all of
\(\mathrm{PW}^{\mathrm{hol}}_{F}\) is not helped by this reduction.
\end{proof}

\begin{remark}[elementary ceiling]
\label{v4-arith:w6-a:rem:three-rule-outs}
No elementary reformulation of \((\star)\) via (i) spectral projection
to the core, (ii) monotone-majorant replacement, or (iii)
Paley--Wiener support shrinkage improves the sufficient condition of
\Cref{v4-arith:w6-a:thm:HFD-ATP}. Every such elementary construction
recovers at most the HFD sub-class.
\end{remark}

\section{The Irreducible Obstruction, Precisely Named}
\label{v4-arith:w6-a:irreducible}

\begin{theorem}[irreducible A-TP residual]
\label{v4-arith:w6-a:thm:irreducible-ATP}
\ClaimStatusProvedHere. The archimedean inequality \(W_{\infty}\ge0\)
unconditionally closes on
\(\mathrm{PW}^{\mathrm{hol},\mathrm{HFD}}_{F}\) (the high-frequency
dominance sub-class of \Cref{v4-arith:w6-a:def:HFD-class}). Its
residual on the complement
\(\mathrm{PW}^{\mathrm{hol}}_{F}\setminus\mathrm{PW}^{\mathrm{hol},\mathrm{HFD}}_{F}\)
reduces to the bounded-interval domination statement
\eqref{v4-arith:w6-a:eq:bounded-reduction-star} on
\([-t_{\ast},t_{\ast}]\). This statement is strictly equivalent to
the Weil positive form of RH only with the \(R[f]\) functional
retained, by
\Cref{v4-arith:w5-a:thm:ATP-equals-weil-positivity}, and by
\Cref{v4-arith:w6-a:prop:rule-out-projection},
\Cref{v4-arith:w6-a:prop:rule-out-majorant}, and
\Cref{v4-arith:w6-a:prop:rule-out-pw-shrinkage} it admits no
further elementary reformulation.
\end{theorem}

\begin{proof}
Combine \Cref{v4-arith:w6-a:thm:HFD-ATP} (unconditional
\(W_{\infty}\)-closure on HFD),
\Cref{v4-arith:w6-a:thm:bounded-reduction} (bounded-interval
reduction with \(R[f]\) retained), and the three ruling-out
propositions.
\end{proof}

\begin{remark}[A-TP: proved, open, and ruled out]
\label{v4-arith:w6-a:rem:honest-form}
The direct elementary obstruction yields the following partition.

\begin{description}
\item[\textbf{Proved.}] (a) Strict monotonicity of \(g\) in \(|t|\)
(\Cref{v4-arith:w6-a:lem:g-monotone}), a new elementary fact about
the archimedean Tate kernel that does not appear in any prior
Vol~IV chapter.
(b) Existence and uniqueness of the sign-change \(t_{\ast}\)
(\Cref{v4-arith:w6-a:thm:t-star}).
(c) The decomposition
\(W_{\infty}(f*\widetilde{f})=(B[f]-A[f])/(2\pi)\)
(\Cref{v4-arith:w6-a:lem:decomposition}).
(d) Unconditional \(W_{\infty}\)-positivity on the explicit sub-class
\(\mathrm{PW}^{\mathrm{hol},\mathrm{HFD}}_{F}\)
(\Cref{v4-arith:w6-a:thm:HFD-ATP}).
(e) Reduction of full A-TP to the inequality
\(B[f]\leq A[f]+2\pi R[f]\), with \(A[f]\) supported on
\([-t_{\ast},t_{\ast}]\) and \(B[f]\) supported on its complement
(\Cref{v4-arith:w6-a:thm:bounded-reduction}).
\item[\textbf{Named open.}] The \(R[f]\)-coupled bounded-interval
domination
statement \eqref{v4-arith:w6-a:eq:bounded-reduction-star} for
\(f\in\mathrm{PW}^{\mathrm{hol}}_{F}\setminus\mathrm{PW}^{\mathrm{hol},\mathrm{HFD}}_{F}\).
This is RH-equivalent by
\Cref{v4-arith:w5-a:thm:ATP-equals-weil-positivity} and
unreachable by the three elementary reductions listed in
\Cref{v4-arith:w6-a:bounded-interval:ruled-out}.
\item[\textbf{Ruled out elementarily.}] (i) Spectral projection to the
core (\Cref{v4-arith:w6-a:prop:rule-out-projection}), (ii) monotone
majorant replacement
(\Cref{v4-arith:w6-a:prop:rule-out-majorant}), (iii) Paley--Wiener
width shrinkage
(\Cref{v4-arith:w6-a:prop:rule-out-pw-shrinkage}).
\end{description}
\end{remark}

\subsection{Non-Elementary Residuals}
\label{v4-arith:w6-a:irreducible:open-obstructions}

\begin{remark}[three non-elementary constructions that remain live]
\label{v4-arith:w6-a:rem:live-constructions}
The \(R[f]\)-coupled bounded-interval inequality
\eqref{v4-arith:w6-a:eq:bounded-reduction-star} is not closed by
elementary analysis. Three non-elementary constructions remain open in
principle.

\begin{enumerate}
\item \emph{Spectral density of non-trivial zeros.} The zero-side
functional equals \(R[f]-W_{\infty}(f*\widetilde{f})\) by
\eqref{v4-arith:w6-a:eq:weil-rearranged}. Replacing it by the positive
ordinate-square sum is the RH-level step. A zero-density obstruction must
control the actual zero divisor and only then specialise to ordinates.
The Riemann--von Mangoldt formula
\(N(T)=(T/2\pi)\log(T/2\pi)-T/2\pi+O(\log T)\) is classical; its
sharpening to an integrable weight against the zero-divisor
test values is pursued in
Chapter~\ref{v4-arith:w5-i:chapter}
(\(R_{\mathrm{analytic}}\)).
\item \emph{Gaussian Weil test functions.} For \(f\) a Gaussian, the
core integral \(A[f]\) and tail integral \(B[f]\) admit closed-form
computations via steepest-descent methods. The positivity question
reduces to an explicit finite-dimensional moment problem on the
Gaussian family. Hardy's 1942 approach to the zero-free region of
\(\zeta\) enters here; Chapter~\ref{v4-arith:w5-i:chapter} pursues
this further.
\item \emph{Hermite--Biehler/de Branges spectral structure.}
The bounded-interval domination statement \((\star)\) admits a de
Branges space reformulation: inside the Hilbert space
\(\mathcal{H}(E_{\xi})\) of
Chapter~\ref{v4-arith:w5-j:VBstar-spectral-positivity}, \((\star)\)
is a positivity condition for the evaluation kernel
\(K(t,t')=(\xi_{\mathbb{R}}(1/2+it)\xi_{\mathbb{R}}(1/2-it')-\xi_{\mathbb{R}}(1/2-it)\xi_{\mathbb{R}}(1/2+it'))/(t-t')\).
Its positivity on \([-t_{\ast},t_{\ast}]\) against
\(|\widehat{f}|^{2}\) is equivalent to the Hermite--Biehler
condition HB\(^{+}\) of
\Cref{v4-arith:w5-j:thm:VBstar-main-direction}.
\end{enumerate}

None of (1)--(3) is elementary. Each is a substantive construction
equivalent in strength to the de Branges or
Connes--Deninger constructions. The direct elementary obstruction establishes
that the bounded-interval term is the irreducible archimedean sign
content after all elementary reductions; the finite-prime and polar
terms remain in \(R[f]\). Further change requires entering one of
these three non-elementary constructions.
\end{remark}

\begin{remark}[comparison to Chapter~\ref{v4-arith:w5-a:chapter}]
\label{v4-arith:w6-a:rem:compare-w5}
Chapter~\ref{v4-arith:w5-a:chapter} left A-TP as the archimedean Tate
distribution, positivity strictly equivalent to RH and irreducibly
open. The direct elementary obstruction (Chapter~\ref{v4-arith:w6-a:ATP-direct-obstruction})
refines this:
\begin{itemize}
\item \emph{Explicit sub-class closure}: \(W_{\infty}\ge0\) holds on
\(\mathrm{PW}^{\mathrm{hol},\mathrm{HFD}}_{F}\), a strictly smaller
but explicit and non-empty class.
\item \emph{Bounded-interval reduction}: the residual A-TP confines
its sign-changing archimedean term to an explicit compact interval
\([-t_{\ast},t_{\ast}]\) with \(t_{\ast}\approx 0.8507\), while \(R[f]\)
retains the polar and finite-prime terms.
\item \emph{Elementary ceiling}: no further elementary reformulation
improves the HFD bound; the ceiling is established by three ruling-out
propositions.
\end{itemize}
The RH-equivalent residual persists. The refinement is precision:
the residual is a bounded-interval archimedean sign inequality coupled
to \(R[f]\), not an unbounded archimedean spectral problem.
\end{remark}

\section{Direct Obstruction Boundary}
\label{v4-arith:w6-a:summary}

\begin{enumerate}
\item The digamma kernel \(g(t)=\mathrm{Re}\,\psi(1/4+it/2)+\log\pi\)
is even, strictly monotone in \(|t|\), with
\(g(0)=\psi(1/4)+\log\pi\approx-3.083\) and
\(g(t)\sim\log(|t|\pi/2)\to+\infty\)
(\Cref{v4-arith:w6-a:lem:g-monotone},
\Cref{v4-arith:w6-a:prop:g-boundary}).
\item \(g\) has a unique positive sign-change
\(t_{\ast}\in(0.84,0.86)\), refined numerically to
\(t_{\ast}\approx 0.8507\)
(\Cref{v4-arith:w6-a:thm:t-star},
\Cref{v4-arith:w6-a:prop:t-star-numerical}).
\item The archimedean Weil contribution decomposes as
\(W_{\infty}(f*\widetilde{f})=(B[f]-A[f])/(2\pi)\) with \(A[f],B[f]\geq 0\)
(\Cref{v4-arith:w6-a:lem:decomposition}).
\item The archimedean term \(W_{\infty}\) is non-negative
unconditionally on the high-frequency dominance sub-class
\(\mathrm{PW}^{\mathrm{hol},\mathrm{HFD}}_{F}\)
(\Cref{v4-arith:w6-a:thm:HFD-ATP},
\Cref{v4-arith:w6-a:cor:HFD-explicit}). This sub-class is explicit,
non-empty (\Cref{v4-arith:w6-a:rem:HFD-nonempty}), and strictly
smaller than \(\mathrm{PW}^{\mathrm{hol}}_{F}\)
(\Cref{v4-arith:w6-a:rem:HFD-strict}).
\item A-TP on the complement reduces to the inequality
\eqref{v4-arith:w6-a:eq:bounded-reduction-star}; its negative
archimedean term is supported on \([-t_{\ast},t_{\ast}]\), while
\(R[f]\) retains the polar and finite-prime terms
(\Cref{v4-arith:w6-a:thm:bounded-reduction}).
\item Three further elementary reductions fail: spectral
projection (\Cref{v4-arith:w6-a:prop:rule-out-projection}), monotone
majorant (\Cref{v4-arith:w6-a:prop:rule-out-majorant}), support
shrinkage (\Cref{v4-arith:w6-a:prop:rule-out-pw-shrinkage}).
\item The irreducible archimedean sign obstruction is precisely named:
the bounded-interval term in the inequality
\(B[f]/(2\pi)\leq A[f]/(2\pi)+R[f]\), exactly equivalent to A-TP only
with the \(R[f]\) functional retained
(\Cref{v4-arith:w6-a:thm:irreducible-ATP}).
\item Three non-elementary constructions remain open: zero-density
refinement, Gaussian reduction, Hermite--Biehler/de Branges
spectral structure (\Cref{v4-arith:w6-a:rem:live-constructions}).
\end{enumerate}

\begin{remark}[Beilinson principle]
\label{v4-arith:w6-a:rem:beilinson-closing}
A smaller unconditional theorem on
\(\mathrm{PW}^{\mathrm{hol},\mathrm{HFD}}_{F}\) is preferred to a
larger conjectural claim on \(\mathrm{PW}^{\mathrm{hol}}_{F}\). The
proved content is elementary, explicit, and computable. The open
archimedean sign content is located on a bounded interval and stated
as a precise inequality coupled to \(R[f]\). The elementary ceiling is
established; no further
elementary advance is available without entering one of the three
non-elementary constructions of \Cref{v4-arith:w6-a:rem:live-constructions}.
\end{remark}
